\chapter{Lezione 3 --- Distribuzioni, regolatore digitale, campionamento e teorema di Shannon}

\section{Richiami dalla lezione precedente}

Riprendiamo le tre rappresentazioni introdotte nella lezione precedente:
\begin{itemize}
    \item \textbf{Serie di Fourier} per segnali periodici di periodo \(T\):
    \[
    f(t)=\sum_{n=-\infty}^{+\infty} C_n e^{j n\omega_0 t},
    \qquad
    \omega_0=\frac{2\pi}{T},
    \]
    con coefficienti
    \[
    C_n=\frac{1}{T}\int_{-T/2}^{T/2} f(t)e^{-j n\omega_0 t}\,dt.
    \]

    \item \textbf{Trasformata di Fourier}:
    \[
    \mathcal{F}\{f(t)\}=F(\omega)=\int_{-\infty}^{+\infty} f(t)e^{-j\omega t}\,dt.
    \]

    \item \textbf{Trasformata di Laplace unilaterale}:
    \[
    \mathcal{L}\{f(t)\}=F(s)=\int_{0^-}^{+\infty} f(t)e^{-st}\,dt.
    \]
\end{itemize}

Questi tre strumenti verranno ora usati per introdurre alcune distribuzioni fondamentali
e per arrivare al problema del campionamento.

\section{Distribuzioni}

Una distribuzione non viene definita come una funzione ``ordinaria'', ma attraverso
le sue proprietà integrali quando viene moltiplicata per una funzione regolare.

In modo intuitivo, se \(g\) è una distribuzione e \(f\) è una funzione regolare,
l'oggetto significativo è
\[
\int_{-\infty}^{+\infty} f(t)\,g(t)\,dt.
\]

Questo numero reale o complesso rappresenta l'azione della distribuzione \(g\)
sulla funzione test \(f\).

\subsection{Distribuzioni fondamentali richiamate a lezione}

Le distribuzioni citate negli appunti sono:
\[
\delta(t),\qquad \delta'(t),\qquad \delta''(t),
\]
e inoltre il gradino unitario
\[
1(t)\equiv H(t),
\]
detto anche \textbf{funzione gradino} o \textbf{funzione di Heaviside}.

\section{Delta di Dirac}

La distribuzione fondamentale è la \textbf{delta di Dirac}, definita dalla proprietà
\[
\int_{-\infty}^{+\infty} f(t)\,\delta(t)\,dt = f(0).
\]

Poiché la delta è concentrata nell'origine, si può anche scrivere
\[
\int_{0^-}^{+\infty} f(t)\,\delta(t)\,dt = f(0).
\]

Dal punto di vista intuitivo:
\begin{itemize}
    \item \(\delta(t)\) è nulla per \(t\neq 0\);
    \item la sua area totale vale \(1\);
    \item il suo effetto ha senso solo sotto segno di integrale.
\end{itemize}

\section{Derivata della delta di Dirac}

Negli appunti compare la derivata della delta. Con la convenzione standard delle distribuzioni,
vale
\[
\int_{-\infty}^{+\infty} f(t)\,\delta'(t)\,dt = -f'(0).
\]

Più in generale, per la derivata \(n\)-esima della delta,
\[
\int_{-\infty}^{+\infty} f(t)\,\delta^{(n)}(t)\,dt
=
(-1)^n f^{(n)}(0).
\]

Questa relazione si ottiene per integrazione per parti in senso distribuzionale.

\section{Gradino di Heaviside}

Il gradino unitario si definisce come
\[
H(t)=1(t)=
\begin{cases}
0, & t<0,\\
1, & t>0.
\end{cases}
\]

Dal punto di vista distribuzionale, il gradino è legato alla delta tramite
\[
H(t)=\int_{-\infty}^{t}\delta(\tau)\,d\tau,
\]
e quindi
\[
\frac{d}{dt}H(t)=\delta(t).
\]

Questo è il collegamento fondamentale tra impulso e gradino.

\section{Tabella elementare di trasformate}

Raccogliamo i segnali richiamati negli appunti e le relative trasformate.

\subsection*{Trasformate di Laplace}

\[
\mathcal{L}\{\delta(t)\}=1,
\]
\[
\mathcal{L}\{H(t)\}=\frac{1}{s},
\]
\[
\mathcal{L}\{tH(t)\}=\frac{1}{s^2},
\]
\[
\mathcal{L}\{e^{-\sigma t}H(t)\}=\frac{1}{s+\sigma},
\qquad \Re(s)>-\sigma.
\]

\subsection*{Trasformate di Fourier}

Nel contesto ingegneristico degli appunti si usano spesso le forme sintetiche:
\[
\mathcal{F}\{\delta(t)\}=1,
\]
\[
\mathcal{F}\{H(t)\}\sim \frac{1}{j\omega},
\]
\[
\mathcal{F}\{tH(t)\}\sim -\frac{1}{\omega^2},
\]
\[
\mathcal{F}\{e^{-\sigma t}H(t)\}=\frac{1}{\sigma+j\omega},
\qquad \sigma>0.
\]

\medskip

\noindent
\textbf{Osservazione.}
In senso pienamente distribuzionale, le trasformate di Fourier di \(H(t)\) e \(tH(t)\)
contengono anche termini proporzionali a \(\delta(\omega)\) e alle sue derivate.
Negli appunti il punto centrale è però la forma operativa usata nei calcoli.

\section{Dal regolatore analogico al regolatore digitale}

Nelle applicazioni reali oggi non si usano quasi più regolatori interamente analogici,
ma regolatori \textbf{digitali}.

Nel caso continuo lo schema classico è:
\[
y^\star \longrightarrow \text{sommatore} \longrightarrow e \longrightarrow R(s) \longrightarrow u \longrightarrow G(s) \longrightarrow y,
\]
con retroazione negativa dell'uscita.

Nel caso digitale, invece, tra errore e impianto compaiono:
\begin{itemize}
    \item un \textbf{ADC} (convertitore analogico-digitale),
    \item un \textbf{microprocessore / PLC} che implementa il regolatore,
    \item un \textbf{DAC} (convertitore digitale-analogico).
\end{itemize}

Quindi il segnale analogico \(e(t)\) viene campionato e trasformato in una sequenza
\[
e[k]=e(kT),
\]
il controllore elabora
\[
u[k],
\]
e il DAC ricostruisce un segnale analogico da applicare all'impianto.

\subsection{Osservazione sul modello dell'impianto}

Negli appunti è sottolineato che il modello dell'impianto deve includere anche
le funzioni di trasferimento dei sensori, se queste sono rilevanti per la dinamica complessiva.

\section{Modello a tempo discreto}

Quando si campiona con periodo \(T\), lo stato viene osservato agli istanti
\[
t=kT, \qquad k\in\mathbb{Z}.
\]

Il modello discreto del sistema assume la forma
\[
x(k+1)=Ax(k)+Bu(k),
\]
dove
\[
x(k)=x(kT).
\]

Questa è la rappresentazione discreta che si usa per il progetto del regolatore digitale.

\section{Operazione di campionamento}

Dato un segnale continuo \(f(t)\), il campionamento consiste nel prelevarne il valore
agli istanti \(t=kT\), ottenendo la sequenza
\[
f[k]=f(kT).
\]

Per descrivere il campionamento nel continuo si introduce il \textbf{pettine di Dirac}
\[
\operatorname{III}_T(t)=\sum_{k=-\infty}^{+\infty}\delta(t-kT).
\]

Il segnale campionato si può allora modellare come
\[
f_s(t)=f(t)\operatorname{III}_T(t)
=
f(t)\sum_{k=-\infty}^{+\infty}\delta(t-kT).
\]

Usando la proprietà della delta si ottiene anche
\[
f_s(t)=\sum_{k=-\infty}^{+\infty} f(kT)\,\delta(t-kT).
\]

Questa scrittura mette in evidenza che il campionamento produce un treno di impulsi
pesati con i valori del segnale agli istanti di campionamento.

\section{Trasformata di Fourier del pettine di Dirac}

Una proprietà fondamentale è che la trasformata di Fourier del pettine di Dirac
è ancora un pettine di Dirac.

Se definiamo
\[
\omega_s=\frac{2\pi}{T},
\]
allora
\[
\mathcal{F}\left\{\operatorname{III}_T(t)\right\}
=
\frac{2\pi}{T}\sum_{n=-\infty}^{+\infty}\delta(\omega-n\omega_s).
\]

Quindi:
\begin{itemize}
    \item nel dominio del tempo il pettine ha passo \(T\);
    \item nel dominio della frequenza il pettine ha passo \(\omega_s=\dfrac{2\pi}{T}\).
\end{itemize}

Se invece si usa la frequenza ordinaria \(\nu=\omega/(2\pi)\), allora la periodicità in frequenza diventa
\[
f_s=\frac{1}{T}.
\]

\section{Digressione sulla convoluzione}

Per capire l'effetto del campionamento in frequenza è utile richiamare la convoluzione.

\subsection{Definizione}

Per due segnali \(f_1(t)\) e \(f_2(t)\), la convoluzione si definisce come
\[
g(t)=f_1(t)\ast f_2(t)=\int_{-\infty}^{+\infty} f_1(\tau)\,f_2(t-\tau)\,d\tau.
\]

Nel caso di segnali causali si può scrivere anche
\[
g(t)=\int_{0}^{t} f_1(\tau)\,f_2(t-\tau)\,d\tau.
\]

La convoluzione è commutativa:
\[
f_1(t)\ast f_2(t)=f_2(t)\ast f_1(t).
\]

\subsection{Uso nei sistemi dinamici}

Se \(g(t)\) è la risposta impulsiva di un sistema lineare tempo-invariante e \(u(t)\)
è l'ingresso, allora l'uscita è
\[
y(t)=u(t)\ast g(t).
\]

Se l'ingresso è la delta di Dirac,
\[
u(t)=\delta(t),
\]
allora
\[
y(t)=g(t).
\]

Questo giustifica il nome di \textbf{risposta impulsiva}.

\subsection{Convoluzione nel dominio di Laplace}

Passando nel dominio di Laplace, la convoluzione diventa prodotto:
\[
\mathcal{L}\{g(t)\}=G(s), \qquad \mathcal{L}\{u(t)\}=U(s),
\]
e quindi
\[
Y(s)=G(s)\,U(s).
\]

\section{Effetto del campionamento nel dominio della frequenza}

Poiché
\[
f_s(t)=f(t)\operatorname{III}_T(t),
\]
nel dominio della frequenza otteniamo una convoluzione:
\[
F_s(\omega)=\frac{1}{2\pi}F(\omega)\ast
\mathcal{F}\{\operatorname{III}_T(t)\}.
\]

Sostituendo la trasformata del pettine:
\[
F_s(\omega)
=
\frac{1}{2\pi}
F(\omega)\ast
\left(
\frac{2\pi}{T}\sum_{n=-\infty}^{+\infty}\delta(\omega-n\omega_s)
\right),
\]
da cui segue
\[
F_s(\omega)
=
\frac{1}{T}\sum_{n=-\infty}^{+\infty} F(\omega-n\omega_s).
\]

Questa è la formula fondamentale del campionamento:
\begin{itemize}
    \item lo spettro del segnale campionato è ottenuto replicando periodicamente
    lo spettro originale;
    \item le repliche sono centrate in multipli della frequenza di campionamento.
\end{itemize}

Nel dominio della frequenza ordinaria \(\nu\), la formula si scrive analogamente come
\[
F_s(\nu)=\frac{1}{T}\sum_{n=-\infty}^{+\infty}F\!\left(\nu-\frac{n}{T}\right).
\]

Quindi lo spettro del segnale campionato è \textbf{periodico in frequenza}
con periodo
\[
\frac{1}{T}.
\]

\section{Interpretazione grafica}

L'effetto geometrico è il seguente:
\begin{itemize}
    \item si parte da uno spettro \(F(\nu)\) centrato attorno a \(0\);
    \item il campionamento genera copie traslate di \(F(\nu)\);
    \item le copie si ripetono attorno a
    \[
    \nu=0,\ \pm \frac{1}{T},\ \pm \frac{2}{T},\dots
    \]
\end{itemize}

Se tali copie non si sovrappongono, il segnale originale può essere ricostruito.
Se invece si sovrappongono, compare aliasing.

\section{Teorema del campionamento}

Supponiamo che \(f(t)\) sia un segnale a banda limitata, cioè che il suo spettro sia nullo
al di fuori di un intervallo
\[
|\nu|\leq \bar{\nu}.
\]

Se la frequenza di campionamento soddisfa
\[
\frac{1}{T}=f_s > 2\bar{\nu},
\]
allora le repliche spettrali non si sovrappongono e il segnale può essere ricostruito
perfettamente con un filtro passa-basso ideale.

La quantità
\[
2\bar{\nu}
\]
è detta \textbf{frequenza di Nyquist} minima richiesta per evitare aliasing.

\subsection{Caso limite}

Se
\[
f_s = 2\bar{\nu},
\]
si è nel \textbf{caso critico} o \textbf{limite di Nyquist}:
le repliche si toccano esattamente ai bordi.

Nella pratica si preferisce lavorare con
\[
f_s > 2\bar{\nu},
\]
in modo da mantenere un margine di sicurezza.

\subsection{Aliasing}

Se invece
\[
f_s < 2\bar{\nu},
\]
le repliche dello spettro si sovrappongono e componenti a frequenze diverse
diventano indistinguibili dopo il campionamento.

Questo fenomeno prende il nome di \textbf{aliasing}.

\section{Significato fisico del teorema}

Il teorema del campionamento dice che un segnale continuo a banda limitata
può essere rappresentato da una sequenza discreta di campioni, purché i campioni
siano abbastanza fitti nel tempo.

In altre parole:
\begin{itemize}
    \item più il segnale contiene frequenze alte, più bisogna campionare velocemente;
    \item se si campiona troppo lentamente, si perde informazione in modo irreversibile.
\end{itemize}

\section{Sintesi della lezione}

I punti chiave della lezione sono:
\begin{enumerate}
    \item richiamo a serie di Fourier, trasformata di Fourier e trasformata di Laplace;
    \item introduzione delle distribuzioni fondamentali:
    \[
    \delta(t),\ \delta'(t),\ \delta''(t),\ H(t);
    \]
    \item relazione tra gradino e delta:
    \[
    \frac{d}{dt}H(t)=\delta(t);
    \]
    \item passaggio dal regolatore analogico al regolatore digitale;
    \item introduzione del modello discreto
    \[
    x(k+1)=Ax(k)+Bu(k);
    \]
    \item modellazione del campionamento tramite pettine di Dirac;
    \item periodicizzazione dello spettro nel dominio della frequenza;
    \item teorema del campionamento e condizione
    \[
    f_s>\,2\bar{\nu};
    \]
    \item comparsa dell'aliasing quando il campionamento è insufficiente.
\end{enumerate}